\documentclass[../main.tex]{subfiles}
\begin{document}

\chapter{Metrics and Evaluation}
\lessoninfo{
  video={Lesson 2, videos 1--25},
  textbook={H\&P \cite{hennessy2017} §1.8--1.9},
  keywords={latency, throughput, speedup, benchmark, geometric mean, iron law, CPI, Amdahl's law, diminishing returns},
}

Before we can improve performance we must be able to measure it, compare
two machines, and predict the effect of a change.  This lesson introduces
the vocabulary (latency, throughput, speedup), the methodology (benchmarks
and how to summarize them) and two analytical tools that recur throughout
the course: the iron law of performance and Amdahl's law.

\section{Performance}

\enquote{Performance} is ambiguous until we say what we are
measuring.\source{Lesson 2, videos 2--3; H\&P §1.8.}

\begin{definition}[Latency]
  The \term{latency}[レイテンシ] of a task is the time from its start to its
  completion.
\end{definition}

\begin{definition}[Throughput]
  The \term{throughput}[スループット] of a system is the number of tasks it
  completes per unit time.
\end{definition}

Throughput is \emph{not} simply the inverse of latency.  When tasks overlap
(pipelining, Lesson 3) or run in parallel (multiprocessors), a system can
complete many tasks per second even though each individual task takes a
long time.\caution{Throughput $=1/\text{latency}$ holds only when tasks are
processed one at a time, strictly one after another.}

\begin{example}
  A processor that takes \SI{10}{\second} for each of two programs run one
  after the other has latency \SI{10}{\second} per program and throughput
  \SI{0.1}{programs/s}.\margin{A DDR4-3200 module returns data
  \SI{14}{\nano\second} after a read command (CL22) yet sustains
  25.6~GB/s, because accesses overlap; Lesson 15.}  A second processor that takes \SI{15}{\second} per
  program but runs two at once has a worse latency yet a throughput of
  $2/15 \approx \SI{0.13}{programs/s}$.
\end{example}

\section{Comparing performance}

To say that machine $X$ is \enquote{$N$ times faster} than machine $Y$ we
use the \term{speedup}[高速化率]:\source{Lesson 2, videos 4--7.}
\begin{equation}
  N \;=\; \frac{\text{speed}(X)}{\text{speed}(Y)}
     \;=\; \frac{\text{latency}(Y)}{\text{latency}(X)}
     \;=\; \frac{\text{throughput}(X)}{\text{throughput}(Y)} .
  \label{eq:speedup}
\end{equation}
Note that latency appears \emph{inverted}: a shorter time means a higher
speed.\caution{A percentage is not a speedup: \SI{30}{\percent} less time is
$N \approx 1.43$, not $1.3$.}  A speedup greater than 1 means $X$ is
better; a speedup below 1 means it is worse (a \enquote{speedup of 0.8} is
a slowdown).  Speedups compose multiplicatively: if $X$ is $2\times$ faster
than $Y$ and $Y$ is $3\times$ faster than $Z$, then $X$ is $6\times$ faster than $Z$.

\tip{Always keep the ratio in the form
$\text{old time}/\text{new time}$; most sign errors in this course come
from inverting it.}

\section{Measuring performance}

The ideal measurement runs the user's actual workload on the actual machine.
This is rarely possible---the machine may not exist yet, and the workload is
usually unknown---so we resort to \term{benchmark}[ベンチマーク]s: programs,
together with their input data, that a community agrees to use as a proxy
for real workloads.\source{Lesson 2, videos 8--10; H\&P §1.8.}  A collection
of such programs is a \emph{benchmark suite}.

\subsection{Types of benchmarks}

From most to least representative:
\begin{description}
  \item[Real applications] complete programs with realistic inputs.  Most
    representative, but hardest to set up, and impossible on a simulator or
    an incomplete prototype.
  \item[Kernels] the most time-consuming part of a real application (e.g.\ a
    matrix multiply), extracted so that it can be run on its own.\margin{A
    cycle-level simulator such as gem5 runs three to four orders of magnitude
    slower than the hardware it models.}  Easier to run; still reflects the
    real work.
  \item[Synthetic benchmarks] programs written to \emph{behave like} real
    applications (similar mix of instructions and memory accesses) while
    being simpler to run.  Useful early in a
    design.\sidenote{The classic examples are Whetstone (1972, floating
    point) and Dhrystone (Weicker, 1984, integer).  Both lost their
    credibility once optimizing compilers began to delete the code whose
    results are never used, which is why a synthetic program is no longer
    accepted as a final measurement.}
  \item[Peak performance] the theoretical maximum rate the hardware can
    sustain (e.g.\ instructions per second if every unit were busy).  A
    marketing number that no real program reaches.
\end{description}

\begin{note}
  Peak performance is not a benchmark result at all; it is an upper bound
  that says nothing about how close real programs get to it.
\end{note}

\subsection{Benchmark standards}

Because a vendor could pick benchmarks that flatter its own product,
independent organizations define standard suites.\source{Lesson 2, video
12; \cite{spec}.}  SPEC (Standard Performance Evaluation Corporation)
publishes CPU, graphics and server suites; TPC does the same for
transaction-processing workloads; EEMBC for embedded systems.  Members of
these organizations agree on the programs, inputs and reporting
rules.\margin{SPEC CPU2017 contains 43 programs in four suites: integer and
floating point, each in a speed and a rate version.  \texttt{spec.org}.}

\section{Summarizing performance}

A suite yields one number per program; we usually want one number per
machine.\source{Lesson 2, videos 13--14; H\&P §1.8.}

\begin{itemize}
  \item \textbf{Execution times} may be combined with the ordinary
    arithmetic mean: the average time is well defined, and the ratio of the
    average times of two machines equals the speedup of the total runtime.
  \item \textbf{Speedups} (ratios) must \emph{not} be averaged
    arithmetically.  The arithmetic mean of ratios depends on which machine
    is the reference, and it does not equal the ratio of the mean times.
    Use the \term{geometric mean}[幾何平均]:
    \begin{equation}
      \overline{S} = \Bigl(\prod_{i=1}^{n} S_i\Bigr)^{1/n} .
      \label{eq:geomean}
    \end{equation}
    \tip{Compute it as $\exp\bigl(\frac{1}{n}\sum_i \ln S_i\bigr)$: a product
    of many ratios overflows, a sum of logarithms does not.}
\end{itemize}

\begin{example}
  Program A takes \SI{10}{\second} on machine $X$ and \SI{20}{\second} on
  $Y$; program B takes \SI{20}{\second} on $X$ and \SI{10}{\second} on
  $Y$.  The per-program speedups of $X$ over $Y$ are 2 and 0.5.  The
  arithmetic mean, 1.25, suggests $X$ is faster; computed with $Y$ as the
  reference the arithmetic mean is also 1.25, suggesting $Y$ is faster---a
  contradiction.  The geometric mean is $\sqrt{2\times 0.5} = 1$ either way,
  correctly reporting that the machines are equivalent on this suite.
\end{example}

\caution{The geometric mean is consistent, not predictive: it does not say
how long the suite takes.  For that, add the times.}

\begin{keypoint}
  Average \emph{times} arithmetically; average \emph{speedups}
  geometrically.  The geometric mean is the only mean for which the mean of
  the ratios equals the ratio of the means.
\end{keypoint}

\section{The iron law of performance}

The \term{iron law}[性能の鉄則] decomposes CPU time into three factors
that map onto the different parts of a computer system:\source{Lesson 2,
videos 15--18; H\&P §1.9.}
\begin{equation}
  \text{CPU time}
  \;=\; \underbrace{\#\text{instructions}}_{\text{IC}}
  \;\times\; \underbrace{\frac{\text{cycles}}{\text{instruction}}}_{\text{CPI}}
  \;\times\; \underbrace{\frac{\text{seconds}}{\text{cycle}}}_{T_{\mathrm{clk}}} .
  \label{eq:iron-law}
\end{equation}

\begin{description}
  \item[Instruction count (IC)] determined by the algorithm, the compiler
    and the instruction set (ISA).
  \item[Cycles per instruction (CPI)] determined by the ISA and by the
    processor's microarchitecture (pipelining, out-of-order execution,
    caches, \ldots).  Most of this course is about lowering CPI.
  \item[Clock cycle time ($T_{\mathrm{clk}}$)] determined by the
    microarchitecture (how much logic sits between two
    latches),\margin{Pentium~4 (Prescott): 31 stages,
    \SI{3.8}{\giga\hertz}; the Core of 2006: 14 stages, under
    \SI{3}{\giga\hertz}, and faster.} the circuit design and the fabrication
    technology.
\end{description}

The value of the law is that it exposes trade-offs: a change that reduces
one factor often increases another.  A deeper pipeline lowers
$T_{\mathrm{clk}}$ but raises CPI; a richer ISA lowers IC but may raise CPI
or $T_{\mathrm{clk}}$.\caution{Never evaluate a change by a single factor.
\enquote{Higher clock frequency} is not \enquote{faster} unless CPI and IC
are unchanged.}

\subsection{Unequal instruction times}

When different instruction classes take different numbers of cycles, the
CPI is a weighted average.\caution{Vendors quote IPC $= 1/\text{CPI}$.
Averaging IPCs over a suite has the same defect as averaging speedups.}
With $n$ classes,
\begin{equation}
  \text{CPU time} \;=\; \Bigl(\sum_{i=1}^{n} \text{IC}_i \times \text{CPI}_i\Bigr)
  \times T_{\mathrm{clk}} ,
  \label{eq:iron-law-classes}
\end{equation}
where $\text{IC}_i$ and $\text{CPI}_i$ are the count and cycles per
instruction of class $i$.

\begin{example}
  A program executes $10^{9}$ instructions on a \SI{2}{\giga\hertz}
  processor: \SI{60}{\percent} are ALU operations (1 cycle),
  \SI{30}{\percent} loads/stores (2 cycles) and \SI{10}{\percent} branches
  (3 cycles).\margin{Realistic order of magnitude: SPEC integer programs
  execute about \SI{25}{\percent} loads, \SI{10}{\percent} stores and
  \SI{15}{\percent} branches.  H\&P Appendix A.}  The average CPI is
  $0.6\cdot 1 + 0.3\cdot 2 + 0.1\cdot 3 = 1.5$, so the CPU time is
  $10^{9}\times 1.5 \times \SI{0.5}{\nano\second} = \SI{0.75}{\second}$.
\end{example}

\section{Amdahl's law}

\term{Amdahl's law}[Amdahl の法則] predicts the overall speedup
obtained by speeding up only \emph{part} of the execution.\source{\cite{amdahl1967}; Lesson 2, videos 19--22; H\&P §1.9.}

\begin{definition}[Amdahl's law]
  If an enhancement speeds up a fraction $f_{\mathrm{enh}}$ of the
  \emph{original} execution time by a factor $s_{\mathrm{enh}}$, the overall
  speedup is
  \begin{equation}
    S_{\mathrm{overall}} \;=\;
    \frac{1}{(1 - f_{\mathrm{enh}}) + \dfrac{f_{\mathrm{enh}}}{s_{\mathrm{enh}}}} .
    \label{eq:amdahl}
  \end{equation}
\end{definition}

\caution{$f_{\mathrm{enh}}$ is a fraction of the \emph{original} time, not
of the time after the change, and not a fraction of the instruction count.}

\begin{example}
  Floating-point instructions account for \SI{40}{\percent} of the run
  time.  A new FP unit makes them $4\times$ faster.  The overall speedup is
  $1/(0.6 + 0.4/4) = 1/0.7 \approx 1.43$.  Making the FP unit infinitely
  fast would only reach $1/0.6 \approx 1.67$.\sidenote{Amdahl's three-page
  paper (AFIPS 1967) used this bound against massively parallel machines: with
  \SI{95}{\percent} of the time parallel, no core count exceeds $20\times$.
  Lesson 18 applies it to cores.}
\end{example}

\subsection{Implications}

Letting $s_{\mathrm{enh}} \to \infty$ in \cref{eq:amdahl} gives the upper
bound $S_{\mathrm{overall}} \le 1/(1-f_{\mathrm{enh}})$.  Two design
principles follow.\source{Lesson 2, videos 21--24.}

\begin{keypoint}[Make the common case fast]
  A modest speedup of a large fraction of the time beats a huge speedup of
  a small fraction.  Optimize what the program actually spends its time on.
\end{keypoint}

The lecture also states the converse, jokingly called \enquote{Lhadma's
law}:\caution{\enquote{Lhadma} is \enquote{Amdahl} spelled backwards.  It is
the lecture's own name for the principle, not standard terminology.}
\emph{do not sacrifice the uncommon case too much} when speeding up the
common one.  If the \SI{10}{\percent} that is left untouched by an
optimization is made $10\times$ slower as a side effect, it now takes as
long as the whole original program, and the overall speedup becomes a
slowdown.

\subsection{Diminishing returns}

Repeatedly improving the same part of the execution yields
\term{diminishing returns}[収穫逓減]: once that part has shrunk, it is no
longer the common case, and the remaining time is dominated by everything
else.  After each improvement the architect should re-measure and ask what
the common case is \emph{now}.

\begin{example}
  Starting from the previous example (FP $=\SI{40}{\percent}$ of time),
  the first $2\times$ FP speedup gives $S = 1/(0.6+0.2) =
  1.25$.\margin{Amdahl's law assumes a fixed problem; if it grows with the
  machine, the serial part shrinks (Gustafson 1988).}  A second
  $2\times$ FP speedup on the new machine, where FP is now only
  $0.2/0.8 = \SI{25}{\percent}$ of the time, gives just
  $1/(0.75+0.125) \approx 1.14$.
\end{example}

\begin{keypoint}[Lesson summary]
  Measure with latency or throughput as appropriate; compare with speedups
  and summarize them with the geometric mean.  The iron law
  (\cref{eq:iron-law}) attributes time to IC, CPI and clock period; Amdahl's
  law (\cref{eq:amdahl}) bounds what any partial improvement can achieve.
\end{keypoint}

\begin{marginfigure}
  \resizebox{\linewidth}{!}{% Overall speedup under Amdahl's law for f_enh = 0.4, as a function of the
% speedup of the enhanced part.  The curve saturates at 1/(1-f) = 1.67.
% Plotted as S-1 so that the origin of the picture is S = 1.
\begin{tikzpicture}[x=2.2mm, y=30mm, font=\scriptsize\sffamily,
                    >={Stealth[length=1.4mm]}]
  % asymptote
  \draw[ln-muted, densely dashed] (0,0.667) -- (17,0.667)
    node[pos=0.52, above, font=\tiny\sffamily, text=ln-muted, inner sep=1pt]
    {$1/(1-f_{\mathrm{enh}})$};
  % axes
  \draw[->] (0,0) -- (17.6,0) node[above left, font=\tiny\sffamily, inner sep=1pt]
    {$s_{\mathrm{enh}}$};
  \draw[->] (0,0) -- (0,0.80) node[right, font=\tiny\sffamily, xshift=-1pt] {$S$};
  \foreach \xv in {4,8,12,16} {
    \draw (\xv,0) -- (\xv,-0.015) node[below, font=\tiny\sffamily] {\xv};
  }
  \foreach \yv/\yl in {0/1.0, 0.25/1.25, 0.5/1.5} {
    \draw (0,\yv) -- (-0.25,\yv) node[left, font=\tiny\sffamily, inner sep=1.5pt] {\yl};
  }
  % S = 1/(0.6 + 0.4/s), plotted as S-1
  \draw[very thick, ln-accent]
    (1,0) -- (1.5,0.154) -- (2,0.25) -- (3,0.364) -- (4,0.429) -- (5,0.471)
    -- (6,0.5) -- (8,0.538) -- (10,0.563) -- (12,0.579) -- (16,0.6);
  \foreach \xv/\yv in {1/0, 2/0.25, 4/0.429, 8/0.538, 16/0.6} {
    \fill[ln-accent] (\xv,\yv) circle (1.1pt);
  }
\end{tikzpicture}
}
  \caption{Overall speedup $S$ for $f_{\mathrm{enh}} = 0.4$ as the enhanced
    part gets faster.  Each further improvement of the same part buys less.}
  \label{fig:l02-amdahl}
\end{marginfigure}

\end{document}
